\section{Qualitative Analysis}
\label{sec:qualitative}
\advisor{the cluster produces only quantitative numbers; the paper also needs a thorough qualitative analysis -- (a) the event representation before/after each corruption, and (b) the per-layer statistical distribution of the membrane/spikes, clean vs corrupted.}

The AUROC tables say \emph{how well} each corruption is detected; this section shows
\emph{what each corruption does}, both at the input and inside the network, so the geometry
of Section~\ref{sec:method} can be seen rather than only measured.

\paragraph{(a) Input event representation, before and after.}
Figure~\ref{fig:gallery} renders a representative Gen1 event histogram (ON events in red, OFF
in blue) for the clean input and each corruption across severities. It makes the taxonomy
visible: \texttt{hot\_pixel} sprinkles persistent isolated firings; \texttt{event\_flood}
fills localized patches with dense activity; \texttt{spatial\_dropout} blanks rectangular
regions; \texttt{event\_rate\_shift} uniformly thins or thickens events; \texttt{polarity\_flip}
swaps the red/blue assignment on a fraction of frames (visually striking yet nearly invisible
to a polarity-symmetric network); and \texttt{temporal\_jitter} reorders time bins (subtle in
a single frame, disruptive in the trajectory). Generated by
\texttt{analysis/viz\_corruptions.py}.

\begin{figure}[t]
  \centering
  \todofig{Insert \texttt{outputs/plots/corruption\_gallery.pdf} -- rows = 6 corruptions,
  columns = [clean, sev 1, 3, 5]. Run \texttt{python analysis/viz\_corruptions.py --sample
  <a real Gen1 frame>}; without \texttt{--sample} it renders a synthetic layout frame.}
  \caption{Event representation before/after each corruption (ON = red, OFF = blue).}
  \label{fig:gallery}
\end{figure}

\paragraph{(b) Per-layer membrane-statistics distributions.}
Figure~\ref{fig:perlayer} shows, for each of the four PLIF layers, the distribution of the
three membrane moments $[\mu,\sigma^2,\kappa]$ over channels and frames, clean versus
corrupted. This is the qualitative companion to the theory: \texttt{hot\_pixel} shifts the
$\mu$ distribution bodily (saturation); \texttt{event\_rate\_shift} scales it; the
contractions narrow $\sigma^2$ \emph{toward} the clean mode (visually explaining the AUROC
$<0.5$ inversion); and \texttt{temporal\_jitter} barely moves the shallow layers but distorts
the deep (L4) distribution, where the larger effective time constant integrates timing. The
panels are produced from the already-extracted $\varphi$ tensors by
\texttt{analyse\_plots.plot\_per\_layer\_distributions} (no GPU re-run needed).

\begin{figure}[t]
  \centering
  \todofig{Insert \texttt{outputs/plots/per\_layer\_dist\_<corruption>\_L5.pdf} (a 4$\times$3
  grid: layers $\times$ moments, clean vs corrupted overlaid). Produced by
  \texttt{python analysis/analyse.py} from \texttt{outputs/phi/}.}
  \caption{Per-layer distribution of membrane moments $[\mu,\sigma^2,\kappa]$, clean vs
  corrupted, for a representative contraction (\texttt{spatial\_dropout}) and dilation
  (\texttt{hot\_pixel}).}
  \label{fig:perlayer}
\end{figure}

Together with the sub-threshold $V(t)$ trajectory plots
(\texttt{analyse\_plots.plot\_all\_trajectories}, Section~\ref{sec:theory}), these give the
reader a mechanistic, visual account of why each corruption is or is not detectable.
